\chapter{SYSTEM ANALYSIS}

\section{Existing System}

The existing bus booking system is largely dependent on manual processes or partially automated online platforms. In traditional systems, passengers are required to visit bus counters physically to inquire about routes, seat availability, timings, and fares. This approach is time-consuming, inefficient, and prone to human errors such as incorrect seat allocation or fare calculation. Even in some online platforms, information is often fragmented across different pages, making navigation difficult for users.

Additionally, many existing systems do not provide real-time updates on seat availability, leading to booking conflicts and poor user experience. Administrative tasks such as updating bus schedules, modifying routes, or managing fares are often handled manually or through outdated interfaces. These limitations result in reduced efficiency, lack of transparency, and inconvenience for both passengers and administrators.

\section{Proposed System}

The proposed system introduces a centralized, web-based bus booking platform that simplifies the ticket booking process for passengers and streamlines administrative operations. Passengers can search for buses based on source, destination, and travel date, view available buses, select seats, enter passenger details, and proceed to checkout through a user-friendly interface. The system ensures that only registered and logged-in users can complete the booking process, thereby enhancing security.

For administrators, the system provides a dedicated admin dashboard with special privileges. Admins can add new buses, edit existing bus details, update routes, timings, boarding and dropping points, and modify fare structures as needed. This centralized control improves efficiency, reduces manual effort, and ensures that all data remains consistent and up to date. Overall, the proposed system enhances usability, accuracy, and reliability compared to the existing system.

\section{Feasibility Study}

\subsection{Technical Feasibility}
The proposed system is technically feasible as it is developed using modern and widely supported technologies such as React.js, HTML, CSS, and JavaScript. These technologies ensure smooth performance, responsiveness, and ease of future enhancements. The system architecture supports scalability and can be easily integrated with a backend or database in the future.

\subsection{Economic Feasibility}
The project is economically feasible since it relies entirely on open-source tools and libraries, eliminating licensing and subscription costs. Development and maintenance costs are minimal, making the system cost-effective for long-term use.

\subsection{Operational Feasibility}
The system is designed to be simple and intuitive, ensuring ease of use for both passengers and administrators. Minimal training is required, and users can interact with the system efficiently without technical expertise.

\section{Functional Requirements}

\begin{itemize}
    \item The system must support user registration and login functionality to ensure secure access.
    \item Users should be able to search buses using source, destination, and date filters.
    \item The system must allow seat selection and entry of passenger details based on the number of seats selected.
    \item Checkout functionality should be restricted to logged-in users only.
    \item The system must include an admin login with special privileges.
    \item Administrators should be able to add, edit, and delete bus details, update schedules, routes, and fare information.
\end{itemize}

\section{Non-Functional Requirements}

The system should have a responsive design to support multiple devices such as desktops, tablets, and mobile phones. It must provide a user-friendly and visually appealing interface to enhance user experience. 

Fast navigation and smooth transitions between pages are essential to reduce user effort. The system should be scalable to allow future enhancements such as backend integration, payment gateways, and database connectivity.